\begin{center}
    \textbf{\large Вариант 4}
\end{center}

\samerowans{\begin{problem}{1}
Вычислите: $3816 - 3640 : (330 - 16 \cdot 19) + 84$
\end{problem}}

\vspace{5pt}

\samerowans{\begin{problem}{2}
Решите уравнение: $(18 \cdot 4 - 450 : b) : 3 = 18$
\end{problem}}

\vspace{5pt}

\samerowans{\begin{problem}{3}
Вычислите: $7\text{ т } 56\text{ кг } - 38\text{ ц } + 244\text{ кг}$. Ответ запишите в т, ц, кг.
\end{problem}}

\samerowans{\begin{problem}{4}
Сложили пять одинаковых чисел, из результата вычли 12, затем полученную разность разделили на 6 и получили 13. Какие числа складывали?
\end{problem}}

\vspace{5pt}

\samerowans{\begin{problem}{5}
Миша два часа делал домашнее задание. Три пятых часа он делал математику, одну шестую часа — русский язык. Все оставшееся время он учил английский. Сколько минут Миша учил английский?
\end{problem}}

\vspace{5pt}

\begin{problem}{6} Из двух пунктов А и В, расстояние между которыми 80 км, одновременно навстречу друг другу выехали два экскаватора. Скорость первого экскаватора, выехавшего из п. А, на $5\text{ км/ч}$ больше скорости второго. За два часа им удалось встретиться и разъехаться на расстояние 50 км. На каком расстоянии от п А будет второй экскаватор через 3 часа 20 минут после своего выезда из В?
\end{problem}

\vspace{5pt}

\noindent\grid{36}{19}

\vspace{5pt}

\samerowans{\begin{problem}{7}
Первый станок может сделать 480 деталей за 8 часов, второй справится с этой задачей в 2 раза быстрее. Сколько часов совместной работы потребуется двум станкам, чтобы сделать 540 деталей?
\end{problem}}

\newpage

\begin{problem}{8} Из прямоугольника с периметром 18 дм вырезали квадрат со стороной 3 см. Найди площадь оставшейся фигуры, если ширина прямоугольника 400 мм.
\end{problem}

\smallskip

\noindent\grid{36}{15}

\vspace{5pt}

\samerowans{\begin{problem}{9}
Таня заплатила за покупку 1 м 40 см шелковой ленты 560 рублей. Надя купила ленты на 50 см больше, чем Таня. Сколько рублей заплатила Надя?
\end{problem}}

\vspace{5pt}

\samerowans{\begin{problem}{10}
Коля и Петя сыграли партию в шахматы, которая закончилась победой одного из них. «Я победил», — сказал Коля. «Коля не проиграл», — сказал Петя. Кто из них выиграл, если известно, что хотя бы кто-то из них сказал неправду?
\end{problem}}

\vspace{5pt}

\begin{problem}{11} Через шесть лет мальчик будет втрое старше, чем он был два года назад. А девочка будет через два года вдвое старше, чем пять лет назад. Кто старше: мальчик или девочка?
\end{problem}

\smallskip

\noindent\grid{36}{20}

\newpage

\begin{center}
    \large{Вариант 4}
\end{center}


\begin{enumerate}
    \item \textbf{Вычислите:} $3816 - 3640 : (330 - 16 \cdot 19) + 84$ \\
    \textit{Решение:}
    \begin{itemize}
        \item $16 \cdot 19 = 304$
        \item $330 - 304 = 26$
        \item $3640 : 26 = 140$
        \item $3816 - 140 + 84 = 3760$
    \end{itemize}
    \textbf{Ответ: 3760}

    \item \textbf{Решите уравнение:} $(18 \cdot 4 - 450 : b) : 3 = 18$ \\
    \textit{Решение:}
    \begin{itemize}
        \item $72 - 450 : b = 54$
        \item $450 : b = 18 \implies b = 25$
    \end{itemize}
    \textbf{Ответ: 25}

    \item \textbf{Вычислите:} 7 т 56 кг $- 38$ ц $+ 244$ кг \\
    \textit{Решение:} $7056 \text{ кг} - 3800 \text{ кг} + 244 \text{ кг} = 3500 \text{ кг} = 3 \text{ т } 5 \text{ ц } 0 \text{ кг}$. \\
    \textbf{Ответ: 3 т 5 ц 0 кг}

    \item \textbf{Задача (пять одинаковых чисел):} \\
    \textit{Решение:} $(5x - 12) : 6 = 13 \implies 5x = 90 \implies x = 18$. \\
    \textbf{Ответ: 18}

    \item \textbf{Домашнее задание Миши:} \\
    \textit{Решение:} 2 часа = 120 мин.
    \begin{itemize}
        \item Математика: 36 мин, Русский: 10 мин.
        \item Английский: $120 - 46 = 74$ мин.
    \end{itemize}
    \textbf{Ответ: 74 минуты}

    \item \textbf{Экскаваторы:} \\
    \textit{Решение:}
    \begin{itemize}
        \item Суммарный путь: $80 + 50 = 130$ км. Скорость сближения: 65 км/ч.
        \item $v_1 = v_2 + 5 \implies v_2 = 30, v_1 = 35$ км/ч.
        \item Путь второго из B за $3\frac{1}{3}$ ч: $30 \cdot 10/3 = 100$ км.
        \item Расстояние от A: $|100 - 80| = 20$ км.
    \end{itemize}
    \textbf{Ответ: 20 км}

    \item \textbf{Станки:} \\
    \textit{Решение:}
    \begin{itemize}
        \item Пр1: 60 дет/ч, Пр2: 120 дет/ч.
        \item Время вместе: $540 / 180 = 3$ часа.
    \end{itemize}
    \textbf{Ответ: 3 часа}

    \item \textbf{Площадь фигуры:} \\
    \textit{Решение:} $P = 180$ см, ширина $b = 40$ см $\implies$ длина $a = 50$ см. \\
    $S = (50 \cdot 40) - (3 \cdot 3) = 1991 \text{ см}^2$. \\
    \textbf{Ответ: 1991 см$^2$ }

    \item \textbf{Шелковая лента:} \\
    \textit{Решение:} 1 см стоит 4 руб. Надя купила 190 см $\implies$ $190 \cdot 4 = 760$ руб. \\
    \textbf{Ответ: 760 рублей}

    \item \textbf{Шахматы (Логика):} \\
    \textit{Решение:} Оба утверждения означают победу Коли. Так как кто-то из них солгал, победил Петя. \\
    \textbf{Ответ: Петя}

    \item \textbf{Возраст:} \\
    \textit{Решение:}
    \begin{itemize}
        \item Мальчик: $x + 6 = 3(x - 2) \implies x = 6$ лет.
        \item Девочка: $y + 2 = 2(y - 5) \implies y = 12$ лет.
    \end{itemize}
    \textbf{Ответ: девочка (мальчику 6 лет, девочке 12 лет)}
\end{enumerate}